\documentclass[10pt,a4paper]{amsart}
\usepackage[margin=19mm]{geometry}
\usepackage{amsmath,amssymb,mathtools}
\usepackage[scr=rsfs,cal=euler]{mathalfa}
\usepackage{xfrac}
\usepackage[unicode,hidelinks,bookmarks=false]{hyperref}
\newcommand{\ie}{i.e.\ }
\newcommand{\cf}{cf.\ }
\newcommand{\resp}{respectively\ }
\newcommand{\Subsection}{\subsection}
\providecommand{\nobreakdash}{\nobreak\hbox{-}}
\setlength{\emergencystretch}{2em}
\multlinegap0pt
\usepackage{mathtools}
\usepackage{lmodern}
\hypersetup{
  colorlinks=true,
  linkcolor=blue,
  citecolor=blue,
  urlcolor=blue
}
\newtheorem{theorem}{Theorem}
\newtheorem{proposition}[theorem]{Proposition}
\newtheorem{lemma}[theorem]{Lemma}
\newtheorem{corollary}[theorem]{Corollary}
\newtheorem{definition}[theorem]{Definition}
\newtheorem{remark}[theorem]{Remark}
\newcommand{\C}{\mathbb C}
\newcommand{\R}{\mathbb R}
\newcommand{\D}{\mathbb D}
\newcommand{\re}{\operatorname{Re}}
\newcommand{\im}{\operatorname{Im}}
\newcommand{\rank}{\operatorname{rank}}
\newcommand{\id}{\operatorname{id}}
\begin{document}
\title[A Lewy-type theorem for pluriharmonic mappings in C^2 and bi]{A Lewy-type theorem for pluriharmonic mappings in $\C^2$ and bi-Lipschitz quasiconformal harmonic maps}
\author{David Kalaj}
\date{}
\maketitle
\noindent\emph{Source and licence.} A Lewy-type theorem for pluriharmonic mappings in $\C^2$ and bi-Lipschitz quasiconformal harmonic maps. Version arXiv:2608.11948v2 (CC BY 4.0). This material is distributed under the Creative Commons Attribution 4.0 International licence, http://creativecommons.org/licenses/by/4.0/, as recorded in the arXiv OAI licence metadata for this exact version and shown on the abstract page https://arxiv.org/abs/2608.11948v2. Full rights notice: licenses/arxiv-2608.11948-CC-BY-4.0.txt.\par\smallskip
\noindent\textit{Editorial scope.} Counted unit: the original \texttt{theorem} environment \texttt{th:local-factorization} at source lines 390--405 together with its complete proof at lines 407--420; the delivered document keeps source lines 278--421 so that every referenced label and every citation resolves inside it. Native editable source is the paper's own concatenated TeX; the AMS wrapper is editorial formatting only. No publisher class, upstream script or unrelated artwork is redistributed.
\section{Analytic preliminaries}

We collect two standard local facts used in the proof of the Lewy-type
theorem: the existence of holomorphic primitives for real-valued
pluriharmonic functions, and the local factorization of plane curve germs.
We begin with the elementary primitive result used in the proof of the main
theorem.

\begin{lemma}\label{lem:primitive}
Let $U\subset\C^2$ be a simply connected domain, and let
$u\in C^2(U,\R)$ be pluriharmonic. Then there is a holomorphic function
$f$ on $U$ such that
\[
        u=\re f.
\]
The function $f$ is unique up to addition of a purely imaginary
constant.
\end{lemma}

\begin{proof}
Since $u$ is pluriharmonic,
\[
        \bar\partial\partial u=0.
\]
Set
\[
        \omega=2\partial u.
\]
Then $\omega$ is a holomorphic $(1,0)$-form. Moreover,
\[
        d\omega=\partial\omega+\bar\partial\omega=0,
\]
because $\partial^2=0$ and $\bar\partial\partial u=0$. Since $U$ is
simply connected, the closed holomorphic one-form $\omega$ has a
holomorphic primitive. Hence there exists a holomorphic function $f$
such that
\[
        df=\omega=2\partial u.
\]
Taking real parts gives
\[
        d(\re f)=\re(df)=\re(2\partial u)=du.
\]
Thus $u-\re f$ is constant. Adding a real constant to $f$, we may
arrange that $u=\re f$. If $f$ and $g$ are two such holomorphic
functions, then $\re(f-g)=0$, so $f-g$ is constant and purely imaginary.
\end{proof}

\subsection{Local factorization of plane curve germs}

We shall use standard facts from local analytic geometry and plane curve
singularities. References include Grauert--Remmert
\cite[Ch.~2]{GrauertRemmert1984}, H\"ormander \cite[Ch.~VI]{Hormander1990},
Chirka \cite{Chirka1989}, Wall \cite[Ch.~2]{Wall2004}, and
Brieskorn--Kn\"orrer \cite[Ch.~8]{BrieskornKnorrer1986}.

\begin{definition}\label{def:convergent-power-series}
The ring
\[
        \C\{z_1,z_2\}
\]
is the local ring of convergent power series at $0\in\C^2$. Thus an
element of $\C\{z_1,z_2\}$ is a power series
\[
        \sum_{j,k\ge0} a_{jk}z_1^jz_2^k
\]
which converges in some neighborhood of the origin, considered as a germ
at $0$.
\end{definition}

\begin{definition}\label{def:unit-irreducible-associated}
Let $R=\C\{z_1,z_2\}$.
\begin{enumerate}
\item A function $u\in R$ is a \emph{unit} if it has a multiplicative
inverse in $R$. Equivalently, $u(0)\ne0$. After shrinking the
representative neighborhood, this means that $u$ is nowhere zero.

\item A nonzero nonunit $g\in R$ is \emph{irreducible} if every
factorization
\[
        g=ab,
        \qquad a,b\in R,
\]
forces either $a$ or $b$ to be a unit.

\item Two nonzero elements $g,h\in R$ are \emph{associated} if
\[
        g=vh
\]
for some unit $v\in R$. Otherwise they are \emph{nonassociated}.
\end{enumerate}
\end{definition}

Geometrically, an irreducible factor in $\C\{z_1,z_2\}$ defines one
local branch of a plane analytic curve. Associated irreducible factors
define the same branch, while pairwise nonassociated irreducible factors
define distinct branches.

\begin{proposition}\label{prop:ufd-convergent-power-series}
The local analytic ring
\[
        \C\{z_1,z_2\}
\]
is a unique factorization domain.
\end{proposition}

\begin{proof}
The ring $\C\{z_1,z_2\}$ is a regular local ring. Every regular local
ring is factorial. Hence $\C\{z_1,z_2\}$ is a unique factorization
domain. See, for example, Matsumura \cite[Theorem~20.3]{Matsumura1986}.
\end{proof}


\begin{theorem}[Local factorization theorem]\label{th:local-factorization}
Let $f$ be a nonzero holomorphic function near $0\in\C^2$, and suppose
that
\[
        f(0)=0.
\]
After shrinking the neighborhood of $0$, one can write
\[
        f=u f_1^{m_1}\cdots f_r^{m_r},
\]
where $u$ is holomorphic and nowhere zero, the functions
$f_1,\ldots,f_r$ are irreducible elements of $\C\{z_1,z_2\}$, no two of
them are associated, and $m_j\ge1$. The germs $\{f_j=0\}$ are precisely
the distinct irreducible local branches of the plane curve germ
$\{f=0\}$, and $m_j$ is the multiplicity of the corresponding branch.
\end{theorem}

\begin{proof}
The Taylor series of $f$ defines an element of $\C\{z_1,z_2\}$. By
Proposition~\ref{prop:ufd-convergent-power-series}, this ring is a
unique factorization domain. Hence $f$ has a factorization
\[
        f=u f_1^{m_1}\cdots f_r^{m_r},
\]
where $u$ is a unit, the $f_j$ are pairwise nonassociated irreducible
elements, and $m_j\ge1$. Since $u(0)\ne0$, after shrinking the
representative neighborhood the function $u$ is nowhere zero. Thus $u$
does not affect the zero set. The irreducible factors $f_j$ give
exactly the distinct local analytic branches of $\{f=0\}$, and the
exponent $m_j$ records the multiplicity of the corresponding branch.
\end{proof}


\begin{thebibliography}{99}
\bibitem[BrieskornKnorrer1986]{BrieskornKnorrer1986} E. Brieskorn and H. Kn\"orrer, \newblock \emph{Plane Algebraic Curves}, \newblock translated by J. Stillwell, \newblock Birkh\"auser, Basel, 1986.
\bibitem[Chirka1989]{Chirka1989} E. M. Chirka, \newblock \emph{Complex Analytic Sets}, \newblock Kluwer Academic Publishers, Dordrecht, 1989.
\bibitem[GrauertRemmert1984]{GrauertRemmert1984} H. Grauert and R. Remmert, \newblock \emph{Coherent Analytic Sheaves}, \newblock Grundlehren der mathematischen Wissenschaften, Vol. 265, \newblock Springer-Verlag, Berlin, 1984.
\bibitem[Hormander1990]{Hormander1990} L. H\"ormander, \newblock \emph{An Introduction to Complex Analysis in Several Variables}, \newblock third revised edition, \newblock North-Holland, Amsterdam, 1990.
\bibitem[Matsumura1986]{Matsumura1986} H. Matsumura, \newblock \emph{Commutative Ring Theory}, \newblock Cambridge Studies in Advanced Mathematics, Vol. 8, \newblock Cambridge University Press, Cambridge, 1986.
\bibitem[Wall2004]{Wall2004} C. T. C. Wall, \newblock \emph{Singular Points of Plane Curves}, \newblock London Mathematical Society Student Texts, Vol. 63, \newblock Cambridge University Press, Cambridge, 2004.
\end{thebibliography}
\end{document}
